% s-093-p093-m002; original PDF 93, 94
\Aphorism{s-093-p093-m002}{98}
Vielleicht giebt es in Fragen höherer Cultur kein schärferes Werthmass für den Einzelnen als die Art seiner \Sperr{Polemik}, wobei ich nicht zu sagen brauche, nach welcher Seite die steigenden, nach welcher die fallenden Werthe zu rechnen sind. Ein durchaus polemischer Mensch ist eine Schande der Natur, ein heilloser Tartuffe, etwas hoffnungslos Unproductives ohne Aufgabe, Sinn und Zweck. Diesen Idealisten ohne Ideal dankt man es, dass unter tüchtigen und positiven Naturen der „Idealist“ als ein grundverlogenes Wesen gilt, das zwischen sein Wirken und Wollen einen Widerspruch hineinkeilt, nicht um vom Einen aus das Andere zu reformiren, sondern um des Einen wie des Anderen überhoben zu sein. Leider lässt sich ihm kein vollkommener Gegentypus gegenüberstellen, da kaum einmal ein höherer Mensch ganz ohne Polemik hat leben können, selbst wenn er ohne polemische Anlage geboren war. Um so entschiedener ist daran festzuhalten, dass jede Polemik zunächst Sünde wider den Geist und nur aus tiefsten NothwehrInstincten des höheren Menschen allenfalls zu rechtfertigen ist; wo Polemik ohne Noth, ohne Herausforderung geübt wird, ist es eben nicht der höhere Mensch, der sie übt. Wer Mühe und Kosten nicht zu scheuen hätte, könnte hiernach durch den Versuch entscheiden, ob ein gegebenes Individuum „Idealist“ oder höherer Mensch ist: er stelle ihm nur die Mittel zur Verfügung, sein Ideal wirklich auszuleben (also in erster Reihe Geld, freie Zeit und Unabhängigkeit, denn auf diese boshaften Trivialitäten kommt heute fast jedes geistige Martyrium hinaus). Der hohe, productive Mensch wird von diesem Tage an seine Polemik einstellen und die Verminderung des äusseren Widerstandes als Glück empfinden; er wird von nun an, nicht buchstäblich, sondern andeutungsweise verstanden, so leben wie Nietzsche im Engadin und am Mittelmeer gelebt hat, \Foreign{procul negotiis} und fern von allem, was „das Inn’re stört“. Für den unproductiven verlogenen ldealisten hingegen ist Widerstand, Störung, \Foreign{profanum vulgus} eine Nothwendigkeit; in ruhigen durchsichtigen Lebensformen müsste sich ja zeigen, mit welchem positiven Lebensinhalt er selber geladen ist, während er so allenfalls das verhinderte Genie weiterspielen kann. Zu der einfachen Consequenz, still bei Seite zu treten und sich ein eigenes Lebens- und Glücksgebiet abzugrenzen, wenn es die Umstände gestatten — dazu hat sich noch kein „Idealist“ verstanden; lieber bleibt er beim Alten, um über das Alte gedeihlich und mit gutem Gewissen zu schimpfen.
\par

